% !PRINTED 97
\weberpage{97}

Setzt man
\[
\Delta_\nu=\begin{vmatrix}
 A_1^{(1)}&\cdots&A_\nu^{(1)}&A_{\nu+1}^{(1)}&\cdots&A_n^{(1)}\\
 \cdots&&\cdots&\cdots&&\cdots\\
 A_1^{(\nu)}&\cdots&A_\nu^{(\nu)}&A_{\nu+1}^{(\nu)}&\cdots&A_n^{(\nu)}\\
 0&\cdots&0&1&\cdots&0\\
 \cdots&&\cdots&\cdots&&\cdots\\
 0&\cdots&0&0&\cdots&1
\end{vmatrix},\tag{5}
\]
wobei $n-\nu$ Zeilen und Colonnen beigefügt sind, von denen die ersteren ausser in den Diagonalgliedern lauter Nullen haben. Bildet man jetzt das Product $A\Delta_\nu$, so folgt
\begin{equation}
A\Delta_\nu=
\begin{vmatrix}
 A&0&\cdots&0&a_{\nu+1}^{(1)}&\cdots&a_n^{(1)}\\
 0&A&\cdots&0&a_{\nu+1}^{(2)}&\cdots&a_n^{(2)}\\
 \cdots&\cdots&&\cdots&\cdots&&\cdots\\
 0&0&\cdots&A&a_{\nu+1}^{(\nu)}&\cdots&a_n^{(\nu)}\\
 0&0&\cdots&0&a_{\nu+1}^{(\nu+1)}&\cdots&a_n^{(\nu+1)}\\
 \cdots&\cdots&&\cdots&\cdots&&\cdots\\
 0&0&\cdots&0&a_{\nu+1}^{(n)}&\cdots&a_n^{(n)}
\end{vmatrix},\tag{6}
\end{equation}
und dies ist nach dem Satz IX, §. 22
\[
= A^\nu
 \begin{vmatrix}
 a_{\nu+1}^{(\nu+1)}&\cdots&a_n^{(\nu+1)}\\
 \cdots&&\cdots\\
 a_{\nu+1}^{(n)}&\cdots&a_n^{(n)}
 \end{vmatrix}.
\]
Dividirt man hier durch $A$ und wendet die Bezeichnung des §. 23 an, so folgt
\begin{equation}
 \Delta_\nu=A^{\nu-1}A_{1,2,\ldots,\nu}^{1,2,\ldots,\nu}.\tag{7}
\end{equation}
Für $\nu=2$ ergiebt sich das specielle Resultat
\begin{equation}
 A_1^{(1)}A_2^{(2)}-A_1^{(2)}A_2^{(1)}=A A_{1,2}^{1,2}.\tag{8}
\end{equation}
Die Formel (8) werden wir später öfter benutzen. In der Bezeichnung durch die Differentialquotienten lässt sie sich so darstellen
\begin{equation}
 A\frac{\partial^2A}{\partial a_1^{(1)}\partial a_2^{(2)}}
 =\frac{\partial A}{\partial a_1^{(1)}}\frac{\partial A}{\partial a_2^{(2)}}
 -\frac{\partial A}{\partial a_1^{(2)}}\frac{\partial A}{\partial a_2^{(1)}}.\tag{9}
\end{equation}
Besonders wichtig ist sie in dem Fall, wo $A$ eine symmetrische Determinante ist, wo also $a_i^{(k)}=a_k^{(i)}$ ist; dann ist auch
\[
 \frac{\partial A}{\partial a_1^{(2)}}=\frac{\partial A}{\partial a_2^{(1)}},
\]
%CONT>>
